\documentclass[11pt,a4paper]{article}
\usepackage[utf8]{inputenc}
\usepackage[margin=1in]{geometry}
\usepackage{amsmath,amssymb,amsthm}
\usepackage{algorithm,algorithmicx,algpseudocode}
\usepackage{graphicx}
\usepackage{hyperref}
\usepackage{booktabs}

\newtheorem{theorem}{Theorem}

\title{High-Dimensional Private Empirical Risk Minimization by Greedy Coordinate Descent}
\author{Pierre Cornilleau, \& Amine Belkasmi}
\date{}

\begin{document}

\maketitle

\begin{abstract}
We summarize the main ideas and results of ``High-Dimensional Private Empirical Risk Minimization by Greedy Coordinate Descent'' by Mangold et al.\ (AISTATS 2023). The paper proposes a differentially private greedy coordinate descent (DP-GCD) method that achieves near-optimal utility in high-dimensional settings by exploiting quasi-sparsity and coordinate-wise smoothness, without explicit constraints or prior sparsity knowledge.
\end{abstract}

\section{Introduction}

We consider empirical risk minimization (ERM) problems of the form
\[
  w^* \in \arg\min_{w \in \mathbb{R}^p} f(w) = \frac{1}{n} \sum_{i=1}^n \ell(w,d_i),
\]
under \((\varepsilon,\delta)\)-differential privacy constraints. In high-dimensional regimes, standard differentially private methods typically suffer utility bounds that grow polynomially in the dimension \(p\), which makes them ineffective when \(p \gg n\).

Existing approaches that achieve logarithmic dependence on \(p\) usually require either \(\ell_1\)-constraints (for instance, Frank--Wolfe on an \(\ell_1\) ball), prior knowledge of sparsity, or access to auxiliary public data. Mangold et al.\ instead design an unconstrained optimization algorithm, \emph{DP-GCD}, that automatically adapts to quasi-sparse structure and achieves utility bounds with only logarithmic dependence on \(p\) in favorable regimes.

\section{The DP-GCD Algorithm}

The starting point is greedy coordinate descent (GCD), which updates a single coordinate at each iteration, chosen according to a greedy rule based on the gradient. DP-GCD privatizes both the choice of coordinate and the update along that coordinate.

\subsection*{Algorithm}

Let \(M_j > 0\) denote coordinate-wise smoothness constants and \(L_j > 0\) coordinate-wise Lipschitz constants for the loss. DP-GCD runs for \(T\) iterations and uses Laplace noise scales \(\sigma_j,\sigma'_j\) and step sizes \(\gamma_j\).

\begin{algorithm}[h]
\caption{DP-GCD (Differentially Private Greedy Coordinate Descent)}
\label{alg:dpgcd}
\begin{algorithmic}[1]
\State \textbf{Input:} initial point \(w_0 \in \mathbb{R}^p\), iterations \(T\), smoothness constants \(\{M_j\}_{j=1}^p\), noise scales \(\{\sigma_j,\sigma'_j\}_{j=1}^p\), step sizes \(\{\gamma_j\}_{j=1}^p\).
\For{$t = 0$ \textbf{to} $T-1$}
  \State Compute full gradient \(g_t = \nabla f(w_t)\).
  \State \textbf{Private greedy selection:}
  \[
    j_t \in \arg\max_{j \in [p]} \left\{ \frac{|g_{t,j}|}{\sqrt{M_j}} + \text{Lap}(\sigma_j) \right\}.
  \]
  \State \textbf{Noisy gradient coordinate:}
    \(\tilde g_{t,j_t} = g_{t,j_t} + \text{Lap}(\sigma'_{j_t})\).
  \State \textbf{Update:}
    \(w_{t+1} = w_t - \gamma_{j_t} \tilde g_{t,j_t} e_{j_t}\).
\EndFor
\State \textbf{Output:} \(w_T\).
\end{algorithmic}
\end{algorithm}

Intuitively, the algorithm spends privacy budget only on the most informative coordinates: the index of the chosen coordinate plus a single noisy gradient value per iteration.

\section{Privacy and Utility Guarantees}

\subsection{Privacy}

Assume that for each coordinate \(j\), the per-example loss is \(L_j\)-Lipschitz in \(w_j\). Then the averaged gradient coordinate has \(\ell_1\)-sensitivity \(\Delta^j_1 = 2L_j/n\). DP-GCD uses Laplace noise both in the Report-Noisy-Max step and in the noisy gradient update.

\begin{theorem}[Privacy, informal]
\label{thm:privacy}
Let \(T \ge 1\) and fix \(\varepsilon,\delta > 0\). Suppose the noise scales satisfy
\[
  \sigma_j = \sigma'_j
  = \Theta\!\left(\frac{L_j}{n} \frac{\sqrt{T \log(1/\delta)}}{\varepsilon}\right)
  \quad \text{for all } j \in [p].
\]
Then Algorithm~\ref{alg:dpgcd} is \((\varepsilon,\delta)\)-differentially private after \(T\) iterations.
\end{theorem}

\paragraph{Proof sketch.}
At each iteration, data is accessed twice: once through the noisy scores used in Report-Noisy-Max, once through the noisy gradient coordinate. Each access is calibrated with Laplace noise according to the sensitivity \(\Delta^j_1 = 2L_j/n\), yielding per-step guarantees \(\varepsilon_t\) determined by \(\sigma_j,\sigma'_j\). With \(2T\) such mechanisms, an advanced composition argument shows that choosing \(\varepsilon_t \approx \varepsilon / \sqrt{T \log(1/\delta)}\) is sufficient to ensure overall \((\varepsilon,\delta)\)-DP. Solving for \(\sigma_j,\sigma'_j\) in terms of \(\varepsilon_t\) gives the stated scaling.

\subsection{Utility: main rates}

The analysis assumes a coordinate-wise smoothness structure and measures distances in a weighted \(\ell_1\) norm. A key parameter is
\[
  R_{M,1} = \max_{w} \|w - w^*\|_{M,1},
\]
which plays the role of an effective radius in that norm. Under convexity and a suitable choice of \(T\), DP-GCD achieves an excess risk bound of the form
\[
  f(w_T) - f^* \;=\; \tilde O\left( \frac{R_{M,1}^{4/3} \, L_{\max}^{2/3} \, \log^{1/3} p}{n^{2/3} \, \varepsilon^{2/3} \, M_{\min}^{1/3}} \right),
\]
where tildes hide logarithmic factors and \(L_{\max}, M_{\min}\) are extremal coordinate-wise constants.

We focus on the strongly convex case, where the main result is more striking.

\begin{theorem}[Strongly convex utility, simplified]
\label{thm:sc}
Assume:
\begin{itemize}
  \item \(f\) is \(M\)-component-smooth: for all \(w,v\),
  \[
    f(w) \le f(v) + \langle \nabla f(v), w-v\rangle
      + \frac{1}{2} \|w-v\|_{M,2}^2,
  \]
  with \(\|x\|_{M,2}^2 = \sum_{j=1}^p M_j x_j^2\).
  \item \(f\) is \(\mu_{M,1}\)-strongly convex with respect to \(\|\cdot\|_{M,1}\), i.e.
  \[
    f(w) \ge f(v) + \langle \nabla f(v), w-v\rangle
      + \frac{\mu_{M,1}}{2} \|w-v\|_{M,1}^2.
  \]
  \item Noise scales satisfy Theorem~\ref{thm:privacy}, and
    \(\gamma_j = 1/M_j\).
\end{itemize}
Then, for a suitable choice of \(T = \Theta(\mu_{M,1}^{-1} \log(\cdot))\),
\[
  f(w_T) - f^* \;=\; \tilde O\left(
    \frac{L_{\max}^2 \, \log^2 p}{\mu_{M,1} \, M_{\min}^2 \, n^2 \, \varepsilon^2}
  \right).
\]
\end{theorem}

\subsection{Sketch of strongly convex analysis}

We outline the main steps behind Theorem~\ref{thm:sc}.

\paragraph{Step 1: noisy descent inequality.}
By \(M\)-component smoothness and the update \(w_{t+1} = w_t - \gamma_{j_t} \tilde g_{t,j_t} e_{j_t}\) with \(\gamma_{j_t} = 1/M_{j_t}\), one obtains
\[
  f(w_{t+1}) \le f(w_t)
    - \frac{1}{M_{j_t}} \nabla_{j_t} f(w_t) \tilde g_{t,j_t}
    + \frac{1}{2M_{j_t}} \tilde g_{t,j_t}^2.
\]
Writing \(\tilde g_{t,j_t} = \nabla_{j_t} f(w_t) + \xi_t\) with Laplace noise \(\xi_t\) and expanding, this leads to a bound of the form
\[
  f(w_{t+1}) \le f(w_t)
    - \frac{1}{2M_{j_t}} |\nabla_{j_t} f(w_t)|^2
    + \frac{1}{2M_{j_t}}\bigl(\xi_t^2 + 2|\nabla_{j_t} f(w_t)|\,|\xi_t|\bigr).
\]

\paragraph{Step 2: greedy selection advantage.}
The (noisy) greedy rule ensures that the chosen coordinate captures a controlled fraction of the gradient energy,
\[
  \frac{1}{M_{j_t}} |\nabla_{j_t} f(w_t)|^2
  \gtrsim \frac{\|\nabla f(w_t)\|^2}{\sum_{j} M_j}
    - \text{(terms depending on }\sigma_j\text{)},
\]
with high probability. Intuitively, if a coordinate had much smaller contribution, some other coordinate would dominate the noisy score and be chosen instead.

\paragraph{Step 3: strong convexity link.}
Strong convexity with respect to \(\|\cdot\|_{M,1}\) and a duality argument imply that
\[
  \|\nabla f(w_t)\|^2 \gtrsim \mu_{M,1} \bigl(f(w_t) - f^*\bigr),
\]
up to constants involving sums of \(M_j\). Combining this with Step 2 shows that the deterministic part of the update contracts the suboptimality by a factor \(\rho < 1\) at each iteration, ignoring noise.

\paragraph{Step 4: recursion and noise accumulation.}
Substituting the bounds from Steps 2--3 into the noisy descent inequality yields
\[
  f(w_{t+1}) - f^* \;\le\; \rho \bigl(f(w_t) - f^*\bigr)
    + C_1\,\xi_t^2 + C_2\,|\nabla_{j_t} f(w_t)|\,|\xi_t|,
\]
for constants \(C_1,C_2\) depending on \(M_{\min},M_{\max}\). Taking expectations and using bounds on the moments of Laplace noise, together with the noise calibration from the privacy analysis, leads to a cleaner recursion
\[
  \mathbb{E}[f(w_{t+1}) - f^*]
  \le \rho\,\mathbb{E}[f(w_t) - f^*]
    + C\,\sigma'^2,
\]
where \(\sigma'^2\) is of order \(L_{\max}^2 T \log(1/\delta)/(n^2\varepsilon^2)\).

Unrolling this recursion for \(T\) steps gives an exponentially decaying term \(\rho^T(f(w_0)-f^*)\) and a noise accumulation term of order \(T \sigma'^2\). Choosing \(T \asymp \mu_{M,1}^{-1}\log(\cdot)\) balances these two contributions and yields the rate in Theorem~\ref{thm:sc}, including the \(\log^2 p\) factor coming from the dependence of \(\rho\) on \(p\).

\subsection{Quasi-sparse structure}

A central feature of DP-GCD is its ability to exploit quasi-sparsity. If the optimal solution \(w^*\) has only \(s\) coordinates significantly larger than zero and the iterates remain \(s\)-sparse (for instance when starting from \(w_0 = 0\)), the bounds can be sharpened so that the dependence on the ambient dimension \(p\) is replaced by a dependence on \(s\):
\[
  f(w_T) - f^* \;=\; \tilde O\left( \frac{L_{\max}^2 \, s^2 \, \log^2 p}{\mu_{M,2} \, M_{\min} \, n^2 \, \varepsilon^2} \right).
\]

In other words, the effective dimension governing the utility becomes \(s\) (up to logarithmic factors), and DP-GCD automatically adapts to this structure without requiring prior knowledge of \(s\).

\section{Experiments}

The experimental section compares DP-GCD with differentially private stochastic gradient descent (DP-SGD) and a private randomized coordinate descent baseline (DP-CD) on both synthetic and real datasets, with varying sample sizes \(n\) and dimensions \(p\).

Synthetic datasets include moderately high-dimensional logistic regression problems (``log1'', ``log2'') where quasi-sparsity is induced through log-normal weights, and a high-dimensional linear regression problem (``square'') with exactly 10 non-zero coefficients in a \(p=1000\)-dimensional space. Real datasets include a low-dimensional regression task (``california'') and very high-dimensional classification tasks such as ``dorothea'' with \(p \approx 88\,000\) features.

Results show that:
\begin{itemize}
  \item On high-dimensional sparse or quasi-sparse problems (``square'', ``dorothea''), DP-GCD is the only method that substantially decreases the objective, while DP-SGD and DP-CD struggle to make progress under the same privacy budget.
  \item On medium-dimensional quasi-sparse problems (``log1'', ``log2'', ``madelon''), DP-GCD converges in fewer passes over the data and achieves lower relative error than the baselines.
  \item On low-dimensional problems without obvious sparsity (``california''), DP-GCD performs comparably to the baselines, and its per-iteration overhead is not clearly justified.
\end{itemize}

Overall, the experiments support the theoretical message: DP-GCD is particularly effective when the dimension is large and the underlying model is (approximately) sparse.

\section{Discussion and Conclusion}

\paragraph{Strengths.}
DP-GCD is an unconstrained method that achieves logarithmic dependence on the ambient dimension in regimes where quasi-sparsity and coordinate-wise regularity hold. It provides near-optimal rates in the strongly convex setting, adapts automatically to effective sparsity without requiring explicit \(\ell_1\) constraints or prior knowledge of the support, and is empirically competitive on high-dimensional benchmarks.

\paragraph{Limitations.}
Each iteration requires a full gradient computation, so the per-iteration cost can be high compared to methods that operate on mini-batches or single coordinates. The main theory focuses on smooth objectives; extensions to non-smooth regularizers (via proximal variants) are proposed but not analyzed in detail. Finally, the method is not uniformly optimal across all regimes: when only \(\ell_2\)-type bounds are available and sparsity is weak, simpler methods such as DP-SGD can be more appropriate.

\paragraph{Takeaway.}
Mangold et al.'s work illustrates how careful mechanism design and exploitation of problem structure can substantially mitigate the curse of dimensionality in private learning. In applications where high dimensionality and approximate sparsity coexist, DP-GCD offers an appealing combination of theoretical guarantees and empirical performance.

\bibliographystyle{plain}
\begin{thebibliography}{1}

\bibitem{mangold2023high}
P.~Mangold, A.~Bellet, J.~Salmon, and M.~Tommasi.
High-dimensional private empirical risk minimization by greedy coordinate descent.
In \emph{AISTATS}, 2023.

\end{thebibliography}

\end{document}
